\stfhIndex{I092}{Fuluobeiniwusipeidui}{弗罗贝尼乌斯配对（Frobenius pairing）}
\section{\personFrobenius{}配对与对偶性}
\label{sec:frobenius-pairing-duality}

回忆~\eqref{eq:20-flat-qh}、\eqref{eq:20-sharp-qh} 中给出的同构
$\flat:QH^*(M)\to QH_*(M)$ 及其逆
$\sharp:QH_*(M)\to QH^*(M)$。仍分别以 $a^\flat,b^\sharp$ 表示这些映射下的像。

\begin{definition}[\personFrobenius{}配对]
\label{def:20-frobenius-pairing}
定义非退化配对
\begin{equation}
 \Delta=\langle\cdot,\cdot\rangle:
 QH^k(M)\otimes QH^{n-k}(M)\longrightarrow\Lambda_\omega^{(0)}
 \label{eq:20-frobenius-pairing}
\end{equation}
为
\begin{equation}
 \left\langle\sum_Aa_Aq^A,\sum_Bb_Bq^B\right\rangle
 =\sum_{C;\,c_1(C)=0}\left(\sum_A(a_A,b_{C-A})\right)q^C,
 \label{eq:20-frobenius-coefficients}
\end{equation}
其中 $(a_A,b_B)$ 是 $H^k(M,\mathbb Q)$ 与 $H_k(M,\mathbb Q)$ 之间的典范\personPoincare{}配对：
\[
 (a_A,b_B)=a_A\cap PD(b_B).
\]
\footnote{译者注：原书把第二因子的次数印作 $n-k$，而本书以 $\dim M=2n$，且本段
随即把 $H^k$ 与 $H_k$ 配对；通常的上同调\personPoincare{}配对应涉及 $2n-k$。原书在
命题 20.4.4 又重复 $n-k$，故译文照录并登记此疑点。}
由 $QH^*(M)$ 定义中的有限性条件，式~\eqref{eq:20-frobenius-coefficients} 的和总是
有限的，因而定义良好。
\end{definition}

\begin{proposition}
\label{prop:20-frobenius-cup}
以 $[M]$ 表示基本类，以 $1=PD[M]$ 表示 $QH^*(M)$ 的单位元。则
\[
 \langle a,b\rangle=(a\cup_Qb)\cap[M]=\langle(a\cup_Qb),1\rangle.
\]
\end{proposition}

\begin{proof}
只需对上同调类 $a,b\in H^*(M)$ 证明第一个等式。这时左端就是经典\personPoincare{}对偶配对
\[
 (a_A,b_{C-A})=a_A\cap PD(b_{C-A})=(a_A\cup b_{C-A})\cap[M].
\]
现计算右端。以闭链 $P_1,P_2$ 代表 $PD(a),PD(b)$。依定义，
$(a\cup_Qb)\cap[M]$ 位于 $H_0(M;\Lambda_\omega)\cong\Lambda_\omega$，由
\[
 \sum_{\alpha\in\pi_2(M)}
 \mathcal M_3(J_0;\alpha)_{(\operatorname{ev}_1,\operatorname{ev}_2,
 \operatorname{ev}_0)}\times(P_1\times P_2\times M)
 T^{\omega(\alpha)}e^{c_1(\alpha)}
\]
给出，其中 $c_1(\alpha)=0$，故
$(a\cup_Qb)\cap[M]\in\Lambda_\omega^{(0)}\subset\Lambda_\omega$。

另一方面，\termFiberProduct{}
$\mathcal M_3(J_0;\alpha)_{(\operatorname{ev}_1,\operatorname{ev}_2,
\operatorname{ev}_0)}\times(P_1\times P_2\times M)$ 定义一个包含在
$C_0(M;\mathbb Q)$ 中的闭链，故维数必须为零。这只有在 $\alpha=0$ 时才可能；
于是 $\mathcal M_3(J_0;\alpha)$ 的所有元素都是常值映射，否则模空间的维数必须至少
有两个实维数，对应第零个标记点 $z_0$ 的变化。注意，由于 $M$ 是整个空间，对
$z_0$ 没有限制，因此可自由移动 $z_0$ 而不影响映射，也不影响在标记点 $z_1,z_2$
处分别与 $P_1,P_2$ 的给定相交。这证明上述\termFiberProduct{}恰好化为经典帽积
$(a\cup_Qb)\cap[M]=(a,b)$。
\end{proof}

定义 $\mathbb Q$-双线性配对 $\Pi$：
\begin{equation}
 \Pi(a,b)=\pi_0\circ\Delta(a,b)=\pi_0(\langle a\cdot b,1\rangle),
 \label{eq:20-rational-pairing}
\end{equation}
其中 $\pi_0:\Lambda_{0,\omega}\to\mathbb Q$ 是显然的投影，
$\pi_0(\sum_Aa_Aq^A)=a_{A=0}$。

\begin{remark}
\label{rem:20-continuous-dual}
即使作为 $\mathbb Q$-向量空间，$QH_*(M)$ 的对偶向量空间 $(QH_*(M))^*$ 也\emph{不同构}于
标准量子上同调 $QH^*(M)$。上述配对实际诱导单射
\[
 QH^*(M)\lhook\joinrel\longrightarrow(QH_*(M))^*,
\]
其像位于 $QH_*(M)$ 上的\emph{连续}线性泛函集合中；这里连续性相对于
式~\eqref{eq:20-downward-valuation} 的\termValuation{} $v$ 所诱导的拓扑而言。
\end{remark}

\begin{proposition}
\label{prop:20-pi-nondegenerate}
$\Pi$ 非退化。换言之，若对所有 $b\in QH^{n-k}(M)$ 都有 $\Pi(a,b)=0$，则 $a=0$。
\end{proposition}

\begin{proof}
若 $\Pi(a,b)=0$，则 $\pi_0(\Delta(a,b))=0$，故
\[
 \sum_{C;\,c_1(C)=\omega(C)=0}\sum_A(a_A,b_{C-A})=0.
\]
但 $c_1(C)=\omega(C)=0$ 蕴含 $C=0$ 于 $\Gamma$，因此对所有
$b_{-A}\in H^{n-k}(M)$ 都有 $\sum_{A\in\Gamma}(a_A,b_{-A})=0$。若
$a\in QH^k(M)$，取 $b=a_{-A}q^{-A}$，其中 $a_{-A}\in QH^{n-k}(M)$，便证明
对所有 $A\in\Gamma$ 都有 $a_A=0$，故 $a=0$。
\footnote{译者注：原书直接取 $b=a_{-A}q^{-A}$，却未说明为何系数 $a_{-A}$ 同时
属于互补次数，亦未取与 $a_A$ 配对非零的\personPoincare{}对偶元素；此处类型与非退化性
论证均有缺口。译文照录。}
\end{proof}

在第 22.5.1 小节中，我们将把 $QH^*(M)$ 看成小\personMorse{}函数的\personFloer{}上同调，利用
$QH^*(M)$ 与 $HF_*(H)$ 之间的自然同构，进一步研究\personFrobenius{}配对
\eqref{eq:20-frobenius-pairing}、\personFloer{}复形配对 (19.7.49) 以及二者的关系。
